\section{Fourier comparison under a leaf law}\label{sec:histograms}

On a leaf with density at most \(2^{cN}\) relative to the raw law,
an event of raw probability \(p\) can have conditional probability as
large as \(2^{cN}p\). To compare the two laws we therefore need an
exponent that can exceed any prescribed \(c\). A single query gives
no such exponent. We estimate a high even moment by sampling many
queries at the same orientation.

On a binary vector space we normalize Fourier coefficients by
\[
 \widehat f(\xi)=\E_x f(x)(-1)^{\xi\cdot x}.
\]
Parseval reads \(\sum_\xi|\widehat f(\xi)|^2=\E|f|^2\).
The Fourier norm is \(\|\widehat f\|_1=\sum_\xi|\widehat f(\xi)|\).
A linear character on a tuple factors into characters of its positions.

\subsection{The leaf estimate}

Fix a query template and its quadratic Gram reference set \(\cK_n\).
Extend a function on \(\cK_n\) by zero to the full key vector space.
Before imposing the unary tests, the nominal input space is affine and
the key map \(K_o(z)\) is affine. Identify that affine space with a binary
vector space by choosing its origin.

% lax begin Lax793645.FourierComparison
\begin{theorem}\label{thm:leaf-fourier}
Fix \(c,b_0\) and \(\epsilon>0\). There is \(\delta>0\), independent
of \(n\), such that the following holds eventually. If
\[
 |D|\le1,\qquad \supp D\subseteq\cK_n,\qquad
 \max_\xi|\widehat D(\xi)|\le\delta,
\]
then every leaf law \(\pi\le2^{cN}\mu^2\) satisfies
\[
 \E_{o\sim\pi}\sup_{\|\widehat w\|_1\le b_0}
 \left|\E_zD(K_o(z))\psi(z)w(z)\right|\le\epsilon.
\]
Here \(\psi\) is a function of the nominal input with \(|\psi|\le1\), fixed before the orientation is drawn. The weight
\(w\) may depend on the orientation. All errors are unnormalized, so
an additional orientation filter of size at most 1 is permitted.
\end{theorem}
% lax end

In particular, \(\psi\) may be any of the nominal binary phases obtained
by expanding the indexed mixer conditions. Their rank can grow with \(n\).
We keep these phases inside the moment calculation. The unary status
weights have bounded Fourier norm by
\Cref{sec:moments}.

Put \(f_o(z)=D(K_o(z))\psi(z)\). For each nominal linear character
\(\alpha\), write
\[
 A_\alpha(o)=\E_z f_o(z)(-1)^{\alpha\cdot z}.
\]
We use an even moment \(\ell=2\lfloor an\rfloor\), with \(a>0\) fixed
and sufficiently small. For every \(H>0\), we will choose \(\delta>0\)
so that, uniformly in \(\alpha\),
\begin{equation}\label{eq:moment-target}
 \E_{\mu^2}|A_\alpha(o)|^\ell\le2^{-HN}
\end{equation}
for all sufficiently large \(n\). The construction parameters remain fixed.
Conditional on \(o\), the \(\ell\) factors in the moment can be
evaluated at independent nominal inputs. Each contributes one copy of
\(D\), but the same frame maps all their columns to the ambient space. The nominal phase and the nominal character
are bounded weights fixed before the orientation is sampled.

\subsection{Keeping dependent columns}

There are at most \(k_0\ell\) formal image columns, for a template
constant \(k_0\). Keep these columns even when their nominal inputs
are dependent. If a relation space has basis \(r_1,\ldots,r_d\), then
\begin{equation}\label{eq:relations-fourier}
 1_{\{Xr_i=0\ \forall i\}}
 =2^{-dN}\sum_{\eta_1,\ldots,\eta_d\in\F^N}
       (-1)^{\sum_i\eta_i\cdot Xr_i}.
\end{equation}
The absolute sum of the coefficients is one. Each summand is a linear
character of the original columns, so it factors over the query batches.
A relation involving many batches costs no factors of \(D\).

Large nominal deficiencies are rare. The fresh \(Z\) parts of the
columns are \(p_s\otimes Z\), whose first selector coordinate is \(Z\).
A specified relation space of dimension \(d\) imposes \(dn\) independent
bit equations on the independent \(Z\) vectors for those columns.
With \(u\) columns there are at most \(2^{du}\) ordered choices for a
basis of a \(d\)-dimensional relation space. Summing the resulting
probability bound over the components gives
\begin{equation}\label{eq:nominal-tail}
 \PP(\text{nominal deficiency at least }d)
 \le C2^{-d(n-C\ell)}.
\end{equation}
We may therefore bound this event separately. For the remaining inputs,
the total nominal deficiency is bounded by a fixed constant.

\subsection{Certifying the exact image kernel}

If a linear map sends a dependent list to ambient columns, it preserves
all the relations in that list. Injectivity of the frame also requires
that no \emph{extra} relation hold. The following identity enforces this
second condition by introducing independent test vectors.

Consider \(u\) nominal columns with kernel \(r\), where \(\dim r=d\).
Choose a surjective matrix \(L:\F^u\to\F^{u-d}\) with kernel \(r\),
and put \(J=L\trans L\). Since \(L\) is onto, the columns of \(J\) span the image of
\(L\trans\), namely \(r^\perp\).
For an ambient column map \(X\), let \(Y\) consist of \(u\) independent
uniform ambient vectors. Nondegeneracy of the ambient polar form gives
\begin{equation}\label{eq:auxiliary-kernel}
 1_{\{\ker X=r\}}
 =2^{u(u-d)}1_{\{r\subseteq\ker X\}}
   \E_Y1_{\{\beta_0(X,Y)=J\}}.
\end{equation}
Indeed, under the kernel inclusion, the image of the map
\(y\mapsto\beta_0(X,y)\) is contained in \(r^\perp\).
It contains every column of \(J\) precisely when \(\ker X=r\).
In that case its rank is \(u-d\), and each prescribed pairing column
has probability \(2^{-(u-d)}\).

Apply \eqref{eq:auxiliary-kernel} in each component, using
\eqref{eq:relations-fourier} for the kernel inclusion. This leaves the
original quadratic Gram equations and the pairings between \(X\) and
\(Y\). The auxiliary vectors are independent and have no prescribed
pairings with one another.

Write \(b(u)=u(u+1)/2\). Equal completion counts identify the image law with the uniform law
conditioned on its exact kernel and quadratic Gram data. The specified
\(d\) independent relations cost \(2^{-dN}\). On the quotient the
remaining \(u-d\) independent columns have prescribed quadratic Gram
data, to which \eqref{eq:quadratic-gram-count} applies. Thus the
conditioning event has probability at least
\[
 \tfrac14\,2^{-dN-b(u-d)}.
\]
Expanding all \(b(u)\) original equations and all \(u^2\) auxiliary
equations, and including the factor in \eqref{eq:auxiliary-kernel}, costs
\[
 2^{-b(u)-u^2+u(u-d)}=2^{-b(u)-ud}.
\]
After normalization the factor is at most
\[
 4\,2^{dN+b(u-d)-b(u)-ud}\le4\,2^{dN}.
\]
The terms quadratic in \(u\) have cancelled. Multiplying over the
fixed number of components costs \(2^{O(dN)}\) when each nominal
deficiency is at most \(d\).

\subsection{One query with its auxiliary vectors}

Assign one auxiliary vector to every original column in its batch.
This doubles the fixed batch size. A permitted quadratic phase in one
batch has an original part and a bilinear part joining original and
auxiliary vectors. No term pairs two auxiliary vectors.

If the bilinear part is nonzero, one of its blocks is the
nondegenerate ambient pairing \(\beta_0\). Averaging the character in
the auxiliary variables forces the corresponding linear combination
of original vectors to equal a prescribed frequency. Cancellation on
this block gives the bound \(2^{-N/4}\) used below.
If the bilinear part is zero, averaging over the auxiliary vectors gives
zero unless their linear frequency vanishes. The original quadratic
phase is a combination of the Gram values fixed on \(\cK_n\), so it is
constant on the support of \(D\). The remaining integral is a Fourier
coefficient of \(D\), of modulus at most \(\delta\).
In either case the local bound is
\begin{equation}\label{eq:auxiliary-local}
 \delta+2^{-N/4}.
\end{equation}
This argument permits arbitrary additional linear characters on the
original and auxiliary columns.

\subsection{Summing the Gram patterns}

Index the Gram coefficient matrix by batches of queries. For disjoint
sets of batches take the submatrix with rows in the first set and
columns in the second. Let \(r\) be the largest rank of any such
submatrix, which we call the cross rank. There are at most
\(2^{C\ell(r+1)}\) coefficient patterns of cross rank \(r\). Indeed,
record the rows and columns meeting an invertible \(r\times r\) minor
that attains this rank, including the other rows and columns in their
batches. This costs \(O(\ell r)\) bits because each batch has bounded
size. Between the remaining distinct batches all entries are determined
by bordering the minor, as in the proof below. The entries within
batches cost a further \(O(\ell)\) bits.

The following lemma handles the patterns of small cross rank. A slot is
the position of a column in its batch. The allowed pairs of slots form
a symmetric relation. All other entries must remain zero.

% lax begin Lax793645.OrderedSlotCompletion
\begin{lemma}[Completing a Gram pattern]\label{lem:ordered-slot-completion}
Let \(M\) be a symmetric binary matrix indexed by batches of \(k\)
slots, supported on the allowed pairs of slots. If its cross rank is
\(r\), there is a decomposition \(M=D+R\) with \(D,R\) symmetric and supported on the
same allowed pairs, \(D\) block diagonal by batch, and
\(\rank R\le2k(1+4k)r\).
\end{lemma}
% lax end
% lax begin Lax793645Proofs.Manuscript.OrderedSlotCompletion.ordered_slot_completion
\begin{proof}
First find \(U\) of rank at most \((1+4k)r\) agreeing with \(M\)
between distinct batches. If \(r=0\), take \(U=0\). Otherwise choose an
invertible \(r\times r\) submatrix \(B=M_{I,J}\) whose row batches
and column batches are disjoint. Start with \(M_{*,J}B^{-1}M_{I,*}\).
For a row and column in distinct batches outside those meeting \(B\),
bordering \(B\) by this row and column gives a singular matrix by
maximality of \(r\). If its added row is \((u,m)\) and its added
column above \(m\) is \(v\), subtracting \(uB^{-1}\) times the old
rows gives \(m-uB^{-1}v=0\). This identifies the entry \(m\).
Restore the rows and columns in the at most \(2r\) batches meeting
\(B\). This adds at most \(4kr\) to the rank.

Order the slots. Let \(V\) keep the entries of \(U\) on allowed slot
pairs \(u\le v\), and let \(W\) keep those on allowed pairs \(u<v\).
Set \(R=V+W^{\mathsf T}\). At distinct slots we copy the same entry
into both positions. At equal slots in different batches, \(U\)
already agrees with the symmetric matrix \(M\). Thus \(R\) is symmetric
and agrees with \(M\) between batches. Each of \(V,W\) is a sum of
\(k\) matrices formed by restricting \(U\) to one row slot and the
corresponding allowed column slots. Both ranks are at most \(k\rank U\).
This gives the asserted bound for \(R\). Take \(D=M-R\).
\end{proof}
% lax end

Fix a rank cutoff \(t\). For \(r\ge t\), the Walsh bound across a rectangle of rank \(r\)
contributes \(2^{-crN}\), even with an additional linear character.
For \(r<t\), apply Lemma~\ref{lem:ordered-slot-completion}. Since \(k\) is
fixed, the correction \(R\) has rank \(O(r)\).
The correction phase uses \(O(r)\) ambient vector combinations and
hence \(O(rN)\) bits. Its character therefore has Fourier norm at most
\(2^{CrN}\). Each term in that Fourier expansion is linear in the
ambient columns and factors by batch. Apply \eqref{eq:auxiliary-local} to every batch. All
\(\ell\) small Fourier factors survive. Together with
\eqref{eq:nominal-tail} and the normalization above, this gives
\begin{align}\label{eq:assembled-moment}
 \E |A_\alpha|^\ell\le{}&C2^{-d(n-C\ell)}\\
 &+2^{C(dN+\ell)}
 \left[2^{CtN+Ct\ell}(\delta+2^{-cN})^\ell
                   +2^{-t(cN-C\ell)}\right].\notag
\end{align}
First choose \(a\) so small that \(n-C\ell\ge n/2\) and
\(cN-C\ell\ge cN/2\). For a prescribed \(H\), increasing \(d\) now
makes the first term less than \(2^{-HN}/3\) for large \(n\).
With \(d\) fixed, increasing \(t\) does the same for the term with
exponent \(-t(cN-C\ell)\), including its prefactor. In the remaining
term all factors outside the power \(\ell\) grow at most as \(2^{C'N}\)
for a fixed \(C'\). Since \(\ell/N\) tends to \(2a/M_0>0\), taking
\(\delta\) sufficiently small makes that term less than \(2^{-HN}/3\)
as well. This proves \eqref{eq:moment-target}.

% lax begin Lax793645Proofs.Manuscript.FourierComparison.one_direct_capped_leaf_fourier
\begin{proof}[Proof of \Cref{thm:leaf-fourier}]
Fourier inversion gives, at each orientation,
\[
 \sup_{\|\widehat w\|_1\le b_0}|\E f_ow|
 \le b_0\max_\alpha|A_\alpha(o)|.
\]
The nominal input space has dimension \(O(n)\), hence at most \(2^{C_0n}\)
linear characters. Markov's inequality and a union bound show that the
leaf mass where the displayed supremum exceeds \(\epsilon/2\) is at most
\[
 2^{cN+C_0n}(2b_0/\epsilon)^\ell
       \sup_\alpha\E_{\mu^2}|A_\alpha|^\ell.
\]
Choose \(H\) in \eqref{eq:moment-target} larger than the fixed exponent
on the right. The exceptional leaf mass is exponentially small. The
supremum is bounded by \(b_0\), so its expectation is eventually at most
\(\epsilon\). The pointwise supremum already allows the weight to depend
on the orientation.
\end{proof}
% lax end

\subsection{Frame bounds for linear cells}

A fixed number of ambient linear equations has a simple counting bound.
Let \(P\) be a raw frame, and let \(P_Z\) be its restriction to the
\(n\) free coordinates. For \(d\) independent ambient rows and any
prescribed matrix \(\Theta\),
\begin{equation}\label{eq:frame-rows}
 \PP_{\rm raw}(LP_Z=\Theta)\le 2^{d^2+2-dn}.
\end{equation}
Here \(d\) is fixed and \(n\) is sufficiently large. To prove this,
identify the rows with ambient vectors by the nondegenerate polar form.
The probability is constant as those vectors vary over their prescribed
quadratic Gram data, by frame transitivity. The Gram count gives this
orbit uniform mass at least \(2^{-d^2-2}\). Average over the orbit and
then drop its defining event. For each fixed injective \(P_Z\), the
pairing map onto its \(n\) coordinates is surjective. Independent
uniform rows therefore have prescribed values with probability
\(2^{-dn}\), proving \eqref{eq:frame-rows}.

Now let \(H\) be an \(r\)-dimensional space of linear forms on a query
with \(k\) ambient vector slots. If its pullback to the free inputs has
kernel dimension at least \(kd\), some slot projection of that kernel
has rank at least \(d\). Choose \(d\) independent rows there and condition
on the other component frames. Their pullbacks under the remaining frame
are prescribed, so \eqref{eq:frame-rows} applies. There are at most
\(2^{kdr}\) maps from \(\F^{kd}\) into \(H\). A union bound gives
\begin{equation}\label{eq:cell-rank-deficit}
 \PP_{\rm raw}(\text{pullback kernel dimension}\ge kd)
 \le 2^{kdr+d^2+2-dn}.
\end{equation}
For a leaf cap \(2^{cN}\), choose \(d>cM_0+11\). For every fixed
\(r\), the exceptional leaf mass is then at most \(2^{-10n}\), eventually.

Outside this exception, averaging translations by the free inputs shows
that each cell has query mass at most \(2^{kd}\) times its ambient
uniform mass. The input laws used here have a fixed density bound
relative to uniform inputs. Include this factor in \(B\) below. The
reference Fourier estimate compares the ambient and Gram cell masses
within a factor of 2. Consequently there is a constant \(B\), chosen
before \(r\), such that
\begin{equation}\label{eq:linear-cell-cap}
 \pi\{o:\PP(K_o\in C)>B\nu(C)\}\le2^{-10n}
\end{equation}
for every cell of every fixed linear partition, once \(n\) is large enough.
